% Part: first-order-logic
% Chapter: syntax-and-semantics
% Section: subformulas

\documentclass[../../../include/open-logic-section]{subfiles}

\begin{document}

\olfileid{fol}{syn}{sbf}

\olsection{\printtoken{P}{subformula}}

\begin{explain}
کسے دتے !!{formula} نوں ”بناؤن“ والے !!{formula}s
دی گل کرنا اکثر فائدہ مند ہوندا اے۔ اسیں اوہناں نوں
اوہدے \emph{!!{subformula}s} آکھدے آں۔ ہر
!!{formula} اپنے آپ دا !!a{subformula} وی اے؛
$!A$ دا اوہدے اپنے آپ $!A$ توں وکھرا ذیلی
فارمولا \emph{پورے توں چھوٹا !!{subformula}} اے۔
\end{explain}

\begin{defn}[براہِ راست !!^{subformula}]
جے $!A$، !!a{formula} اے، تاں $!A$ دے
\emph{براہِ راست !!{subformula}s} اینج استقرائی
طور تے تعریف کیتے جاندے نیں:
\begin{enumerate}
\item ایٹمی !!{formula}s دے کوئی براہِ راست !!{subformula}s نہیں۔

\tagitem{prvNot}{\indcase{!A}{\lnot !B}{$\indfrm$ دا اکلوتا براہِ راست
    !!{subformula}،~$!B$ اے۔}}{}

\item \indcase{!A}{(!B \ast !C)}{$\indfrm$ دے براہِ راست !!{subformula}s
  $!B$ تے $!C$ نیں ($\ast$ دو-مقامی رابطیاں وچوں کوئی وی اک اے)۔}

\tagitem{prvAll}{\indcase{!A}{\lforall[x][!B]}{$\indfrm$ دا اکلوتا براہِ راست
    !!{subformula}،~$!B$ اے۔}}{}

\tagitem{prvEx}{\indcase{!A}{\lexists[x][!B]}{$\indfrm$ دا اکلوتا براہِ راست
    !!{subformula}،~$!B$ اے۔}}{}
\end{enumerate}
\end{defn}

\begin{defn}[پورے توں چھوٹا !!^{subformula}]
جے $!A$، !!a{formula} اے، تاں $!A$ دے
\emph{پورے توں چھوٹے !!{subformula}s} اینج بازگشتی
طور تے تعریف کیتے جاندے نیں:
\begin{enumerate}
\item ایٹمی !!{formula}s دے کوئی پورے توں چھوٹے !!{subformula}s نہیں۔

\tagitem{prvNot}{\indcase{!A}{\lnot !B}{$\indfrm$ دے پورے توں چھوٹے
    !!{subformula}s،~$!B$ تے~$!B$ دے سارے پورے توں
    چھوٹے !!{subformula}s نیں۔}}{}

\item \indcase{!A}{(!B \ast !C)}{$\indfrm$ دے پورے توں چھوٹے !!{subformula}s
  $!B$، $!C$، نالے $!B$ دے سارے پورے توں چھوٹے !!{subformula}s
  تے~$!C$ دے سارے ایہو جہے ذیلی فارمولے نیں۔}

\tagitem{prvAll}{\indcase{!A}{\lforall[x][!B]}{$\indfrm$ دے پورے توں
    چھوٹے !!{subformula}s،~$!B$ تے~$!B$ دے سارے پورے توں
    چھوٹے !!{subformula}s نیں۔}}{}

\tagitem{prvEx}{\indcase{!A}{\lexists[x][!B]}{$\indfrm$ دے پورے توں
    چھوٹے !!{subformula}s،~$!B$ تے~$!B$ دے سارے پورے توں
    چھوٹے !!{subformula}s نیں۔}}{}
\end{enumerate}
\end{defn}

\begin{defn}[!!^{subformula}]
$!A$ دے !!{subformula}s وچ $!A$ خود تے اوہدے
سارے پورے توں چھوٹے !!{subformula}s شامل نیں۔
\end{defn}

\begin{explain}
دھیان دیو کہ اسیں براہِ راست !!{subformula}s تے پورے
توں چھوٹے !!{subformula}s دی تعریف دو ذرا وکھرے
طریقیاں نال کیتی اے۔ پہلی صورت وچ اسیں فارمولا~$!A$
دے ہر ممکن روپ~$!A$ لئی اوہدے براہِ راست !!{subformula}s
دی سیدھی تعریف دتی اے۔ ایہ صورتاں دے لحاظ نال
صاف تعریف اے، تے صورتاں !!{formula}s دے سیٹ
دی استقرائی تعریف دے مطابق نیں۔ دوجی صورت وچ
اسیں سارے !!{formula}s دے سیٹ دی تعریف والا
ڈھنگ فیر اپنایا اے، پر ہر صورت وچ چھوٹے
!!{formula}s $!B$، $!C$ دے پورے توں چھوٹے
!!{subformula}s وی انہاں !!{formula}s دے نال شامل
کیتے نیں۔ ایس نال تعریف \emph{بازگشتی}
بن جاندی اے۔ عام طور تے استقرائی طور تے تعریف
کیتے سیٹ (ساڈے معاملے وچ، !!{formula}s) اُتے
کسے فنکشن دی تعریف بازگشتی اے جے فنکشن دی
تعریف دیاں صورتاں وچ خود اوہ فنکشن ورتیا جائے۔
پر تعریف ٹھیک ہون لئی سانوں یقین کرنا پیندا اے
کہ فنکشن دیاں قدراں صرف اوہناں دلیلاں لئی ورتیاں
جان جیہڑیاں زیرِ تعریف دلیل توں ”پہلاں“ آؤندیاں نیں:
ساڈے معاملے وچ، جدوں $(!B \ast !C)$ لئی
”پورے توں چھوٹے !!{subformula}“ تعریف کردے آں،
تاں صرف ”پہلاں“ والے !!{formula}s $!B$ تے $!C$
دے پورے توں چھوٹے !!{subformula}s ورتدے آں۔
\end{explain}

\begin{prop}
\ollabel{prop:subfrm-trans}
فرض کرو $!B$، $!A$ دا ذیلی فارمولا اے تے
$!C$، $!B$ دا ذیلی فارمولا اے۔ تاں $!C$، $!A$
دا ذیلی فارمولا اے۔ دوجے لفظاں وچ ذیلی فارمولا
تعلق متعدی اے۔
\end{prop}

\begin{prob}
\olref[fol][syn][sbf]{prop:subfrm-trans} ثابت کرو۔
\end{prob}

\begin{prop}
\ollabel{prop:count-subfrms}
فرض کرو $!A$ اک فارمولا اے جیہدے وچ $n$ رابطے
تے مسوّر نیں۔ تاں $!A$ دے زیادہ توں زیادہ $2n+1$
ذیلی فارمولے نیں۔
\end{prop}

\begin{prob}
\olref[fol][syn][sbf]{prop:count-subfrms} ثابت کرو۔
\end{prob}

\end{document}
